\documentclass{article}

\usepackage[english]{babel}

\usepackage[letterpaper,top=2cm,bottom=2cm,left=3cm,right=3cm,marginparwidth=1.75cm]{geometry}

% Useful packages
\usepackage{amsmath}
\usepackage{graphicx}
\usepackage[colorlinks=true, allcolors=blue]{hyperref}
\usepackage{amsthm}
\usepackage{csquotes}
\usepackage{algorithm, algpseudocode, float}


\theoremstyle{definition}
\newtheorem{definition}{Definition}[section]


\title{ Computer Vision Peer Review }
\author{
    Max Moir \\
    \textit{University of Canterbury} \\ 
    \text{Christchurch, New Zealand} \\ 
    \texttt{maxwmoir@gmail.com}
}

\begin{document}
\maketitle

\section{Summary}
% The summary grounds the remainder of your review. You need to demonstrate that you have read and understood the manuscript, which helps the authors understand what other readers are understanding to be the manuscript’s main claims. This is also an opportunity to demonstrate your own expertise and critical thinking, which makes a positive impression on the editors who often may be important people in your field.

% It is helpful to use the following guidelines:

%     Start with a one-sentence description of the paper’s main point, followed by several sentences summarizing specific important findings that lead to the paper’s logical conclusion.
%     Then, highlight the significance of the important findings that were shown in the paper.
%     Conclude with the reviewer’s overall opinion of what the manuscript does and does not do well.


This paper outlines a proposal for the development of a general purpose robot to assist in the caretaking of elderly and disabled people.

The paper then provides preliminary examples related to the area.
It outlines the technical approach that may be taken to build an example robot to meet its criteria.
It then states how said robot may be evaluated and tested.
Ethics / Future work.

It cites two papers that detail applications similar to the intended goal.




\section{Decision}
The final decision of this review is to reject this paper proposal due to the concerns stated below.
The paper requires major revision before it will be in a publishable state.
It requires significant grammatical improvement, and requires conducting new experiments.
Despite the execution falling short in multiple regards, the topic is well motivated.

\section{Major Concerns}


% Speculative 

A first major concern with this paper is that it is speculative rather than experimental.
No clear research contributions.

% Lack of clear experimental or computational data
It follows naturally from this that the paper has a distinct lack of experimental data.

In addition to this, the paper is speculative rather than experimental.
It is clearly stated that the goal of the paper is ``to demonstrate the capabilities of a mobile personal assistance robot using `simple’ image processing.'' 
Despite this, no outcome seems to have been achived in the paper, with the evaluation section being seemingly purely speculative .

% Lack of technical depth



\section{Minor Concerns}

% Intro
In addition to the major concerns about this paper stated above, there are several minor concerns that require revisions.
The paper has many syntactic and grammatical mistakes.

% Poorly worded 

"This paper"

% Grammar

Missing words

    " This a broad topic with many different applications in both health care and non-commercial. "
    items an elderly has

Spelling mistakes:
    (partivuarly repetitive and simple ones)
    Furthermore, logistic needs such as medicine retrieval can be performed without personnel
    hospital personell

incorrect plurals
    follow a similar criteria

missing articles
    Arduino mega 2560 is used for the large number of digital pins

redundancy:
    The robot must be seamless and operate effectively within it’s environment even to the extent of
    operating effectively within an unfamiliar environment that it was not designed for in a way
    that does not compromise human safety

This paper presents a mobile robot for elderly and disabled assistance

missing punctuation
 No obstacle detection is performed, hence the robot will bump into any objects The robot will stop if it cannot find the marker.


corrupted 
    used.
    –¿ arco principes of opereation

% Figure 1: Follower robot not present

% Subsections formatted strangely e.g. Approach
% Paragraph length


% Text on pdf



% Incorrect refence usage:
    oming 5 years.[2] 
    Reference not used [1]
    compromise the user’s safety. []
    [] illustrates a robot 
    With the rise in technology1 the
    However, [2] raises some interesting ethical issues regarding robots in daily lives and health care.
I assume bad citations



\end{document}
